\begin{center}
    {\Huge\bfseries LTE}
\end{center}

\problem{1}{%
Для каждого натурального $n$ докажите существование натурального числа
$m>1$, для которого
\[
  7^n\mid 3^m+5^m-1.
\]}

% ИСТОЧНИК 2 (полностью): 10th Balkan Mathematical Olympiad,
% Nicosia, Cyprus, 1993, Problem 4 (proposed by Romania).
% Полный комплект задач IMO Compendium Group:
% https://imomath.com/othercomp/B/Bmo1993.pdf
% Альтернативная веб-версия с решениями:
% https://schoolexercisebooks.blogspot.com/2010/09/10th-balkan-mathematical-olympiad.html
\problem{2}{%
Пусть $p$ --- простое число, а $m\ge2$ --- целое число. Докажите, что
уравнение
\[
  \frac{x^p+y^p}{2}=\left(\frac{x+y}{2}\right)^m
\]
имеет решение в натуральных числах $(x,y)\ne(1,1)$ тогда и только тогда,
когда $m=p$.}

% ИСТОЧНИК 3 (полностью): Bulgarian Mathematical Olympiad, 1997.
% Формулировка и решение: A. H. Parvardi, "Lifting The Exponent Lemma",
% challenge problem "Bulgaria 1997":
% https://services.artofproblemsolving.com/download.php?id=YXR0YWNobWVudHMvZC83LzY1MGJmN2I1OGU1Y2VmOTQ0NTIwNmExY2FmNDIxY2IxOTdlN2FhLnBkZg%3D%3D&rn=TFRFLnBkZg%3D%3D
% Независимый опубликованный источник с той же атрибуцией и решением:
% https://ibn.idsi.md/sites/default/files/imag_file/86-90_12.pdf
\problem{3}{%
Для некоторого натурального $n$ число $3^n-2^n$ является степенью простого
числа, то есть
\[
  3^n-2^n=q^k
\]
для некоторого простого $q$ и натурального $k$. Докажите, что $n$ --- простое
число.}

% ИСТОЧНИК 4 (полностью): China IMO Team Selection Test 2009, number theory
% problem: a>b>1, b odd, b^n | a^n-1; prove a^b>3^n/n.
% Полный архив китайских математических олимпиад:
% https://upload.xuanyan.ws/chinese-mathematical-olympiads.pdf
% Перепечатка с формулировкой и решением (Problem 8, China TST 2009):
% https://ibn.idsi.md/sites/default/files/imag_file/86-90_12.pdf
\problem{4}{%
Пусть $a>b>1$ --- натуральные числа, причём $b$ нечётно. Известно, что
для некоторого натурального $n$
\[
  b^n\mid a^n-1.
\]
Докажите неравенство
\[
  a^b>\frac{3^n}{n}.
\]}

% ИСТОЧНИК 5 (полностью): IMO Shortlist 1991, Problem 18.
% Полный PDF шорт-листа (задача 18 на странице 3):
% https://www.allameamini.org/wp-content/uploads/2019/03/International_Competitions-IMO_Shortlist-1991-17.pdf
% Дублирующий источник AoPS, A. H. Parvardi, Problem 13:
% https://services.artofproblemsolving.com/download.php?id=YXR0YWNobWVudHMvZC83LzY1MGJmN2I1OGU1Y2VmOTQ0NTIwNmExY2FmNDIxY2IxOTdlN2FhLnBkZg%3D%3D&rn=TFRFLnBkZg%3D%3D
\problem{5}{%
Найдите наибольшее натуральное $k$, для которого $1991^k$ делит число
\[
  1990^{\,1991^{1992}}+1992^{\,1991^{1990}}.
\]}

% ИСТОЧНИК 6 (полностью): 31st International Mathematical Olympiad,
% Beijing, 1990, Day 1, Problem 3.
% Официальный PDF IMO на английском языке:
% https://www.imo-official.org/assets/documents/problems/1990/1990_eng.pdf
% Официальная страница комплекта 1990 года:
% https://www.imo-official.org/problems/1990/
\problem{6}{%
Найдите все целые числа $n>1$, для которых
\[
  n^2\mid 2^n+1.
\]}

% ИСТОЧНИК 7 (полностью): 40th International Mathematical Olympiad,
% Bucharest, 1999, Day 2, Problem 4.
% Официальная страница IMO и PDF комплекта:
% https://www.imo-official.org/problems/1999/
% https://www.imo-official.org/assets/documents/problems/1999/1999_eng.pdf
% Карточка AoPS с точной формулировкой и официальным решением:
% https://artofproblemsolving.com/wiki/index.php?title=1999_IMO_Problems%2FProblem_4
\problem{7}{%
Найдите все пары натуральных чисел $(n,p)$, для которых $p$ --- простое,
$n\le2p$ и
\[
  n^{p-1}\mid (p-1)^n+1.
\]}

% ИСТОЧНИК 8 (полностью): 66th International Mathematical Olympiad,
% Sunshine Coast, Australia, 2025, Day 1, Problem 3 (proposed by Colombia).
% Официальная страница IMO с полной формулировкой:
% https://www.imo-official.org/problems/2025/
% Карточка AoPS задачи:
% https://artofproblemsolving.com/wiki/index.php/2025_IMO_Problems/Problem_3
% Развёрнутые решения и атрибуция автора:
% https://web.evanchen.cc/exams/IMO-2025-notes.pdf
\problem{8}{%
Обозначим через $\mathbb N$ множество натуральных чисел. Функция
$f\colon\mathbb N\to\mathbb N$ называется \textit{суяновой}, если
\[
  f(a)\mid b^a-f(b)^{f(a)}
\]
для любых натуральных $a$ и $b$. Найдите наименьшую действительную константу
$c$, для которой
\[
  f(n)\le cn
\]}
для любой суяновой функции $f$ и любого натурального $n$.